\documentclass[10pt,a4paper]{amsart}
\usepackage[margin=19mm]{geometry}
\usepackage{amsmath,amssymb,mathtools}
\usepackage[scr=rsfs,cal=euler]{mathalfa}
\usepackage{xfrac}
\usepackage[unicode,hidelinks,bookmarks=false]{hyperref}
\newcommand{\ie}{i.e.\ }
\newcommand{\cf}{cf.\ }
\newcommand{\resp}{respectively\ }
\newcommand{\Subsection}{\subsection}
\providecommand{\nobreakdash}{\nobreak\hbox{-}}
\setlength{\emergencystretch}{2em}
\multlinegap0pt
\usepackage{enumerate}
\usepackage{amssymb}
\usepackage[normalem]{ulem}
\usepackage{xcolor}
\newtheorem{theorem}{Theorem}
\title{Balanced intersection size distributions in finite geometries and vector spaces}
\author{Zolt\'an L\'or\'ant Nagy, Zsuzsa Weiner}
\date{Source: arXiv:2605.23644v2 (CC BY 4.0)}
\begin{document}
\maketitle

\noindent\emph{Source and licence.} Balanced intersection size distributions in finite geometries and vector spaces. Version arXiv:2605.23644v2 (CC BY 4.0), distributed by arXiv under the Creative Commons Attribution 4.0 International licence, http://creativecommons.org/licenses/by/4.0/, as recorded in the arXiv licence metadata for this exact version and shown on the abstract page https://arxiv.org/abs/2605.23644v2. Full licence notice: \texttt{licenses/arxiv-2605.23644-CC-BY-4.0.txt}.\par\smallskip

\noindent\textit{Editorial scope.} Counted unit: the article's theorem thm:growing-upper (source0.tex lines 600--617). The article's complete original proof of it is delivered with it (source0.tex lines 619--713, one \texttt{proof} environment of 95 lines). No mathematics is written, repaired or re-derived: the statement and the proof are the article's own lines, in the article's own notation, and no retained line is substituted or re-flowed.  Retained uncounted as the source-local context the proof needs: lines 452--469.  Nothing was omitted from any retained span, so the package carries no omission record.  No retained line contains a citation, so the wrapper carries no bibliography and no bibliography entry.  No retained line contains a figure environment, an image-inclusion command or drawing code, so no image is delivered and the package declares no asset dependency.  Editorial formatting only, in addition: an AMS \texttt{amsart} preamble that restates the article's own declarations and macros used by the retained lines, namely the environments \texttt{theorem} on separate counters, together with the standard packages those lines need.  The retained environments print in this document as Theorem 1 in order of appearance, whereas the article prints the same items with the numbering of its own full document.  The complete native source of the article as distributed in the arXiv e-print for this exact version is bundled unchanged as \texttt{sources/201226/source0.tex}.
	\begin{theorem}
		\label{thm:fixed-upper}
		\label{thm:general-hypergraph}
		Let \(m\) be fixed, and let \((V_n,\mathcal H_n)\) be a sequence of
		\(m\)-uniform block systems, with \(N_n=|\mathcal H_n|\). Suppose that
		\[
		D_n:=|\{(H,H')\in\mathcal H_n^2:H\cap H'\neq\emptyset\}|
		=o(N_n^2)
		\qquad\text{as }n\to\infty.
		\]
		Then
		\[
		B(\mathcal H_n)\leq (\rho_m+o(1))N_n,
		\qquad
		\rho_m=\inf_{0\leq p\leq1}\max_{0\leq i\leq m}
		\binom mi p^i(1-p)^{m-i}.
		\]
	\end{theorem}

	\begin{theorem}
		\label{thm:growing-upper}
		{Let \((V_n,\mathcal H_n)\) be a sequence of \(m_n\)-uniform block systems,
			with \(N_n=|\mathcal H_n|\). Suppose that \(m_n\to\infty\),
			\[
			\max_{\substack{H,H'\in\mathcal H_n\\H\neq H'}}|H\cap H'|=O(1)
			\]
			and
			\[
			m_n\sqrt{\log m_n}=o(N_n).
			\]
			Then
			\[
			B(\mathcal H_n)\leq
			\left(\sqrt{\frac{2}{\pi}}+o(1)\right)
			\frac{N_n}{\sqrt{m_n}}.
			\]}
	\end{theorem}

	\begin{proof} For the upper bound, choose \(S\subseteq V_n\) by including each point independently with probability \(1/2\). The proof follows the same strategy as that of Theorem~\ref{thm:fixed-upper}, but the growing block size requires a sharper covariance estimate in the central range and a separate treatment of the binomial tails. For every \(0\leq i\leq m_n\), a fixed block meets \(S\) in exactly \(i\) points with probability \(\binom{m_n}{i}2^{-m_n}\). Therefore,
		\[
		\mathbb E\bigl[|\mathcal H_n(S,i)|\bigr]
		=N_n\binom{m_n}{i}2^{-m_n}.
		\]
		Moreover,
		\[
		|\mathcal H_n(S,i)|
		=\sum_{H\in\mathcal H_n}\mathbf 1_{\{|H\cap S|=i\}},
		\]
		so its variance is the sum of the covariances of these indicators.
		
		We first consider the central values satisfying
		\(\left|i-m_n/2\right|\leq\sqrt{m_n\log m_n}\). Let \(H\) and \(H'\)
		be distinct blocks and write \(t=|H\cap H'|\). If exactly \(j\) points
		of \(H\cap H'\) belong to \(S\), then the conditional probability that
		\(H\) meets \(S\) in \(i\) points is
		\(\binom{m_n-t}{i-j}2^{-(m_n-t)}\). Since \(t=O(1)\),
		\(0\leq j\leq t\), and \(i\) lies in the central range,
		\[
		\frac{\binom{m_n-t}{i-j}2^{-(m_n-t)}}
		{\binom{m_n}{i}2^{-m_n}}
		=
		2^t\frac{\binom{m_n-t}{i-j}}{\binom{m_n}{i}}
		=
		1+O\left(\sqrt{\frac{\log m_n}{m_n}}\right),
		\]
		where the last estimate follows by writing
		\(i=m_n/2+O(\sqrt{m_n\log m_n})\) and expressing the ratio of the
		binomial coefficients as a product of \(t=O(1)\) factors. Once the
		choices in \(H\cap H'\) are fixed, the two events are independent.
		Their covariance is therefore the variance of the conditional probability
		above. Its relative deviation from
		\(\binom{m_n}{i}2^{-m_n}\) is
		\(O(\sqrt{\log m_n/m_n})\), and hence, uniformly for distinct \(H,H'\),
		\[
		\operatorname{Cov}\left(
		\mathbf 1_{\{|H\cap S|=i\}},
		\mathbf 1_{\{|H'\cap S|=i\}}
		\right)
		=
		O\left(
		\left(\binom{m_n}{i}2^{-m_n}\right)^2
		\frac{\log m_n}{m_n}
		\right)
		=
		O\left(\frac{\log m_n}{m_n^2}\right),
		\]
		where we used
		\(\max_i\binom{m_n}{i}2^{-m_n}=O(m_n^{-1/2})\). The diagonal terms in the variance sum contribute at most \(N_n\binom{m_n}{i}2^{-m_n}=O(N_n/\sqrt{m_n})\), while there are at most \(N_n^2\) off-diagonal terms. Hence
		\[
		\operatorname{Var}\bigl(|\mathcal H_n(S,i)|\bigr)
		=
		O\left(
		\frac{N_n}{\sqrt{m_n}}
		+\frac{N_n^2\log m_n}{m_n^2}
		\right)
		\]
		uniformly in the central range.
		
		There are \(O(\sqrt{m_n\log m_n})\) central values of \(i\). Chebyshev's inequality bounds the probability of a large deviation for each such \(i\), and the union bound adds these probabilities. The variance estimate above and the assumption \(m_n\sqrt{\log m_n}=o(N_n)\) therefore show that, with probability tending to \(1\), simultaneously for every central \(i\),
		\[
		|\mathcal H_n(S,i)|
		=
		N_n\binom{m_n}{i}2^{-m_n}
		+o\left(\frac{N_n}{\sqrt{m_n}}\right).
		\]
		
		It remains to consider the values outside the central range. For every block \(H\), the binomial tail bound gives
		\[
		\Pr\left(
		\left||H\cap S|-\frac{m_n}{2}\right|
		>\sqrt{m_n\log m_n}
		\right)
		\leq 2m_n^{-2}.
		\]
		Thus the expected total number of blocks whose intersection size lies outside this range is at most \(2N_n/m_n^2\). Markov's inequality shows that, with probability tending to \(1\), this number is at most \(N_n/m_n=o(N_n/\sqrt{m_n})\).
		
		Combining the central and tail estimates, there exists \(S\subseteq V_n\) such that
		\[
		M_{\mathcal H_n}(S)
		\leq
		N_n\max_{0\leq i\leq m_n}
		\left\{\binom{m_n}{i}2^{-m_n}\right\}
		+o\left(\frac{N_n}{\sqrt{m_n}}\right).
		\]
		Finally, Stirling's formula gives
		\[
		\max_{0\leq i\leq m_n}
		\left\{\binom{m_n}{i}2^{-m_n}\right\}
		=
		\left(\sqrt{\frac{2}{\pi}}+o(1)\right)m_n^{-1/2},
		\]
		which proves the result.
	\end{proof}

\end{document}
